\chapter{Trabalhos Relacionados}\label{chap:trabalhos_relacionados}

Este capítulo organiza os trabalhos relacionados em três frentes. A primeira reúne sistemas que usam agentes e modelos de linguagem no apoio à revisão científica. A segunda trata das bases de dados que sustentam esses sistemas e da cobertura que elas oferecem. A terceira discute limitações e riscos já documentados no uso de modelos de linguagem como avaliadores.

\section{Revisão científica assistida por agentes e LLMs}

Investigações recentes exploram o papel de agentes baseados em LLMs no processo de revisão por pares. O \textit{framework} Agent Laboratory propõe uma arquitetura de agentes capazes de conduzir diferentes fases do processo de pesquisa, incluindo revisão de literatura, planejamento experimental e redação de relatórios. Os resultados indicam que a colaboração entre humanos e agentes pode reduzir custos e acelerar o desenvolvimento científico, embora o julgamento humano permaneça essencial para garantir a qualidade das avaliações \cite{AgentLaboratory}. O foco desse trabalho, no entanto, é conduzir a pesquisa como um todo, e não avaliar a adequação de um manuscrito a um alvo editorial, que é o foco desta proposta.

Outra linha de pesquisa concentra-se na análise do próprio processo de revisão por pares. O \textit{framework} AgentReview utiliza agentes baseados em LLMs para simular interações entre autores, revisores e editores, permitindo investigar fatores sociológicos e organizacionais que influenciam decisões editoriais. Os experimentos demonstram que vieses individuais e dinâmicas de grupo podem alterar significativamente os resultados de avaliação, evidenciando a complexidade sociotécnica do processo de revisão científica \cite{AgentReview}. Essa evidência apoia uma premissa deste trabalho, a de que a decisão editorial não depende apenas do mérito absoluto do texto, mas o objetivo do AgentReview é estudar o processo, e não orientar o autor antes da submissão.

Há também esforços voltados especificamente à geração automática de pareceres científicos. O método MARG \nomenclature{MARG}{Multi-Agent Review Generation} propõe o uso de múltiplos agentes especializados que analisam diferentes aspectos do artigo, como experimentos, clareza textual e impacto científico. Estudos empíricos mostram que a abordagem multiagente produz comentários mais específicos e úteis do que sistemas baseados em um único agente, indicando vantagens da especialização e da cooperação entre agentes na tarefa de revisão \cite{MARG}. Na mesma direção, o \textit{framework} DeepReview propõe uma abordagem em múltiplos estágios que emula o processo de revisores especialistas, incorporando verificação de novidade por meio de recuperação de literatura, avaliação multidimensional e verificação de confiabilidade baseada em evidências. O modelo resultante, DeepReviewer-14B, foi treinado sobre um conjunto de dados curado com anotações estruturadas e demonstrou desempenho superior ao de sistemas maiores em tarefas de predição de notas e geração de pareceres \cite{DeepReview}.

Esta proposta aproveita duas contribuições desse conjunto de trabalhos. Da arquitetura do MARG, aproveita a ideia de agentes especializados com papéis distintos. Do DeepReview, aproveita o uso de revisões humanas reais como insumo da avaliação. A proposta se distancia deles, porém, em dois pontos. O MARG e o DeepReview avaliam o manuscrito contra critérios de qualidade tomados como universais, como solidez metodológica, novidade e contribuição, enquanto o artefato proposto avalia a adequação a um gênero específico, o de uma trilha de conferência. Além disso, a verificação de novidade do DeepReview, apoiada em recuperação de literatura, introduz um julgamento de mérito absoluto. Se esse julgamento fosse acoplado à verificação de gênero, não seria possível atribuir à modelagem de gênero o efeito observado na avaliação do artefato. Por isso ele não integra o componente avaliado nesta fase, sendo tratado como fronteira de rejeição em trabalho futuro.

\section{Bases de dados de revisão e cobertura}

Os sistemas descritos dependem de conjuntos de dados de revisões humanas em volume, e é aqui que se revela a lacuna que motiva esta pesquisa. As revisões que alimentam esses trabalhos vêm em sua maioria de conferências internacionais de aprendizado de máquina e processamento de linguagem natural, cujos pareceres são disponibilizados pela plataforma OpenReview. O conjunto PeerRead foi um dos primeiros a reunir revisões e decisões editoriais de venues como ACL, NIPS e ICLR para fins de pesquisa \cite{PeerRead2018}. Mais recentemente, o ReviewArena reuniu, em escala, revisões de sete famílias de conferências internacionais, com decisões, réplicas dos autores e o texto integral dos artigos, e constitui hoje uma das maiores bases públicas voltadas à revisão por LLMs \cite{ReviewArena}.

Apesar da escala, essas bases apresentam dois problemas para o objeto desta pesquisa. O primeiro é de idioma e domínio. Elas cobrem apenas conferências internacionais de aprendizado de máquina e processamento de linguagem natural, sem representação de Sistemas de Informação, de comunicação científica em português ou de trilhas temáticas comparáveis às de conferências brasileiras. O segundo é de cobertura da rejeição. Mesmo as bases que preservam pareceres de artigos rejeitados o fazem de modo enviesado. No PeerRead, a seção do NIPS contém revisões apenas de artigos aceitos, e os artigos do arXiv não têm revisão associada, o que limita seu uso para estudar a fronteira entre aceitação e rejeição. Vale notar ainda que a estrutura dos pareceres varia entre conferências e edições nessas bases, o que reforça, em vez de contradizer, a premissa de que o gênero do parecer é específico de cada comunidade.

Não há, até onde se levantou, conjunto público de revisões humanas de artigos científicos brasileiros que inclua artigos rejeitados. Os periódicos nacionais que adotaram avaliação aberta publicam somente pareceres de artigos aceitos, e a Sociedade Brasileira de Computação disponibiliza apenas os anais das conferências, sem os pareceres \cite{solsbc}. Essa ausência tem consequência direta para o projeto. Ela impede, por ora, fundamentar empiricamente um componente que modele a fronteira entre aceitação e rejeição no contexto brasileiro. Por isso o artefato é concebido em duas fases. A primeira modela a conformidade de gênero a partir do corpus de artigos aceitos e pode ser avaliada desde já. 
A segunda dependeria da construção de um corpus de artigos rejeitados, na ausência de um corpus de revisões brasileiras o ReviewArena viabiliza a contrução do segundo agente, pois fornece as revisões negativas reais. E como trabalho futuro, a busca por um corpus de revisões humanas de artigos científicos brasileiros.


\section{Limitações e riscos de LLMs como avaliadores}

O uso de LLMs na avaliação de manuscritos não está livre de falhas, e conhecer essas falhas orienta decisões de projeto do artefato. Li et al. \citeonline{WhereLLMWrong} propõem um \textit{framework} de perturbação multinível guiado por aspectos para diagnosticar as fraquezas de LLMs na revisão automatizada. Por meio de perturbações direcionadas a artigos, revisões e réplicas, os autores identificam vulnerabilidades recorrentes, como a classificação incorreta de falhas metodológicas como problemas de apresentação, a influência desproporcional de recomendações negativas fortes e a interpretação equivocada de críticas incorretas como avaliações rigorosas. Essas fragilidades persistem em diversos modelos amplamente utilizados, o que sustenta a decisão de restringir o artefato ao diagnóstico de conformidade e de manter o juízo final com o autor.

O risco não está apenas no modelo, mas também na forma como o usuário lida com ele. Em um estudo controlado sobre assistentes baseados em LLMs em outro domínio, o da busca de ajuda para uso de software, Khurana, Subramonyam e Chilana \citeonline{WhyandWhen} observaram que os participantes seguiam com frequência as orientações do assistente sem analisá-las, adotando respostas que continham imprecisões e atribuindo as falhas ao próprio desconhecimento, em vez de reconhecer as limitações do modelo. O contexto é diferente do da revisão científica, mas o achado se aplica ao público desta pesquisa. Autores em início de carreira, que ainda não dominam as convenções da comunidade, são justamente os mais propensos a acatar orientações incorretas sem questioná-las. Isso reforça a necessidade de o artefato explicitar o fundamento de cada diagnóstico, ou seja, a qual convenção de gênero ele se refere, para que a orientação possa ser interpretada e contestada pelo autor, e não recebida como veredito.

Em conjunto, esses estudos indicam que agentes de IA podem apoiar a produção científica, mas que sua utilidade depende de delimitar bem o que avaliam e de tornar transparente o motivo de cada apontamento. A lacuna que este trabalho aborda está na interseção de dois recortes ausentes na literatura examinada: o nível da trilha como unidade de análise, tratada como um subgênero com repertório próprio, e o contexto das conferências brasileiras, para o qual não há dados de revisão disponíveis publicamente.